% capitolo 1: introduzione
\chapter{Introduzione}
\label{cap:introduzione}

\section{Contesto macroeconomico}
In ambito macroeconomico, una \textbf{guerra commerciale} è una situazione di conflitto economico tra due o più Paesi, caratterizzata dall'adozione di misure protezionistiche, 
quali l'imposizione di dazi e altre restrizioni agli scambi, in risposta alle politiche commerciali adottate dalle controparti. Tali misure possono innescare meccanismi 
di ritorsione reciproca, alterando i flussi commerciali internazionali e alimentando l'incertezza economica. 
L'imprevedibilità di tali misure genera shock di incertezza che si ripercuotono sui mercati finanziari, contribuendo ad accrescere la volatilità dei rendimenti azionari.

In questo contesto, la guerra commerciale tra Stati Uniti e Cina nel biennio 2025-2026 ha rappresentato un caso emblematico di \emph{escalation} protezionistica, caratterizzato 
da una rapida successione di misure tariffarie e contromisure.

L'\emph{escalation} ha avuto inizio con l'avvio del secondo mandato del Presidente statunitense Donald Trump, quando l'amministrazione americana ha introdotto dazi mirati 
nei confronti di Canada, Messico e Cina, innescando una progressiva intensificazione delle tensioni commerciali e una crescente instabilità sui mercati internazionali.

Le tensioni hanno raggiunto una prima fase di massima intensità nella primavera del 2025, quando, nel mese di aprile, l'aliquota statunitense sulle importazioni cinesi ha 
raggiunto il 125\%. All'\emph{escalation} si sono anche alternati episodi di distensione, come gli accordi di maggio e ottobre 2025.

La vicenda si è successivamente estesa al piano giuridico e istituzionale, culminando nel febbraio 2026 con la pronuncia della Corte Suprema degli Stati Uniti in merito 
all'utilizzo dell'\emph{International Emergency Economic Powers Act} (IEEPA).

Un elemento caratterizzante del periodo è stata la modalità di comunicazione delle decisioni politiche: molte delle decisioni più significative non sono state annunciate 
tramite canali diplomatici o conferenze programmate, bensì attraverso ordini esecutivi inaspettati, annunci sulle piattaforme social e dichiarazioni pubbliche estemporanee, 
spesso rilasciate durante gli orari di chiusura delle principali borse mondiali.

Quando un annuncio viene diffuso a mercati chiusi o durante il fine settimana, il suo effetto non è osservabile subito: i prezzi si aggiustano alla prima seduta utile 
e le diverse \gls{asset-class} possono incorporare la notizia in momenti diversi, a seconda dei fusi orari e degli orari di negoziazione. La data dell'annuncio non coincide quindi 
necessariamente con quella in cui se ne osserva l'effetto e ciò rende necessario selezionare correttamente il momento in cui l'evento entra nell'analisi.

L'analisi assume che gli annunci tariffari siano eventi esogeni decisi indipendentemente dall'andamento dei prezzi, così che le variazioni di volatilità osservate possano 
essere attribuite all'annuncio. Un modello \gls{garch} standard non distingue le variazioni di volatilità dovute a eventi esterni dalle fluttuazioni ordinarie e non permette 
quindi di misurare l'effetto degli annunci.

\section{Obiettivo e metodologie}
Alla luce di queste considerazioni, l'elaborato si pone l'obiettivo di capire se gli annunci tariffari abbiano modificato la volatilità condizionata degli strumenti 
considerati, con quale intensità e persistenza e con quali differenze tra asset class. 
L'analisi considera gli eventi che riguardano i rapporti bilaterali tra Stati Uniti e Cina. A tal fine, l'analisi empirica considera sei strumenti esposti alle tensioni 
commerciali attraverso canali di trasmissione diversi: un indice di mercato (\textbf{S\&P 500}), titoli azionari appartenenti a settori chiave (\textbf{Apple}, \textbf{Nucor} e \textbf{General Motors}), 
l'\textbf{oro} e il mercato valutario, rappresentato dal tasso di cambio \gls{offshore} \textbf{USD/CNH}, caratterizzato da una maggiore flessibilità rispetto al tasso di cambio domestico.

Dal punto di vista metodologico, l'analisi utilizza modelli della famiglia GARCH con componente asimmetrica, nello specifico \gls{gjr-garch} ed \gls{egarch}, con errori distribuiti 
secondo una t di Student. Per stimare l'effetto degli annunci politici, vengono introdotte nell'equazione della varianza le variabili \emph{dummy}, definite su una finestra 
temporale fissa di due sedute per ciascun evento analizzato, con un criterio uniforme per tutti gli eventi.


\section{Organizzazione del testo}
La presente tesi è articolata nei seguenti capitoli: 
\paragraph{Capitolo 1}
Introduce al contesto macroeconomico e agli obiettivi dell'analisi, con una descrizione dei principali strumenti finanziari considerati e delle metodologie utilizzate.

\paragraph{Capitolo 2}
Introduce ai concetti teorici alla base dell'analisi delle serie storiche finanziarie, illustrando le principali metriche di rendimento e i cosiddetti «fatti stilizzati» della volatilità, con particolare attenzione al \emph{volatility clustering} e all'effetto asimmetrico (\emph{leverage effect}). 

\paragraph{Capitolo 3}
Presenta il quadro econometrico di riferimento: i modelli \gls{arch} e GARCH, le loro estensioni asimmetriche (GJR-GARCH ed EGARCH) e la \emph{News Impact Curve}, utilizzata per confrontare graficamente la risposta della volatilità a shock di segno e dimensione diversi.

\paragraph{Capitolo 4}
Descrive dati e metodologia: gli strumenti finanziari selezionati, la cronologia degli annunci tariffari del biennio 2025-2026, la specificazione dei modelli con le variabili \emph{dummy}, le statistiche descrittive e l'analisi esplorativa delle serie. 

\paragraph{Capitolo 5}
Presenta e discute i risultati delle stime, valutando l'impatto degli eventi esogeni sulla volatilità condizionata e la persistenza dei loro effetti. Inoltre, attraverso un'analisi delle correlazioni \emph{rolling}, esamina i co-movimenti tra le diverse asset class nei periodi di maggiore tensione commerciale.

\paragraph{Capitolo 6}
Conclude il lavoro con una sintesi dei risultati ottenuti, con un'analisi critica dei limiti della ricerca e delle possibili estensioni future.


Riguardo la stesura del testo, relativamente al documento sono state adottate le seguenti convenzioni tipografiche:
\begin{itemize}
	\item gli acronimi, le abbreviazioni e i termini ambigui o di uso non comune menzionati vengono definiti nel glossario, situato alla fine del presente documento;
	\item per la prima occorrenza dei termini riportati nel glossario viene utilizzata la seguente nomenclatura: \emph{parola}\glsfirstoccur;
	\item i termini in lingua straniera o facenti parti del gergo tecnico sono evidenziati con il carattere \emph{corsivo}.
\end{itemize}